\documentclass[10pt,twocolumn,a4paper]{article}
% Edycja: 30.10.2026. Kompilacja: pdfLaTeX + BibTeX.
\usepackage[T1]{fontenc}
\usepackage[utf8]{inputenc}
\usepackage{polski,lmodern}
\usepackage[margin=18mm,columnsep=7mm]{geometry}
\usepackage{microtype,graphicx}
\usepackage{booktabs,tabularx,array}
\usepackage{amsmath,amssymb}
\usepackage{enumitem,placeins,xcolor,xurl,float,needspace,etoolbox}
\usepackage[hidelinks,breaklinks=true]{hyperref}
\hypersetup{pdftitle={MOW II -- szablon raportu projektowego}}
\setlist[itemize]{leftmargin=*,itemsep=1pt,topsep=3pt}
\setlist[description]{leftmargin=0pt,labelsep=4pt,itemsep=3pt}
\setlength{\emergencystretch}{1em}
\raggedbottom
\pretocmd{\section}{\Needspace{6\baselineskip}}{}{}
\pretocmd{\subsection}{\Needspace{5\baselineskip}}{}{}
\newcolumntype{Y}{>{\raggedright\arraybackslash}X}
\renewcommand{\abstractname}{Streszczenie}
\renewcommand{\refname}{Bibliografia}

% Dane do zmiany przez zespół.
\newcommand{\TematProjektu}{Tytuł projektu}
\newcommand{\TypZadania}{Klasyfikacja albo regresja}
% Szare akapity są wskazówkami, nie treścią gotowego raportu.
% Po napisaniu własnej treści: \wskazowkifalse
\newif\ifwskazowki
\wskazowkitrue
\newcommand{\wskazowka}[1]{\ifwskazowki
  \par\begingroup\small\color{black!72}\raggedright#1\par\endgroup\fi}
\newcommand{\pole}[1]{\textit{[#1]}}

\title{\TematProjektu\\[2mm]
\large Zaawansowane metody uczenia maszynowego}
\author{Imię i nazwisko autora 1\\\small Numer albumu
\and Imię i nazwisko autora 2\\\small Numer albumu}
\date{\quad \today}
\begin{document}
\maketitle

\begin{abstract}
\wskazowka{W wersji końcowej: problem, dane, protokół oceny,
najważniejszy wynik liczbowy względem baseline i ograniczenie.
W R1 przedstaw cel i plan; zaznacz, że finalne wyniki jeszcze nie istnieją.
Streszczenie aktualizujemy po każdej iteracji, generalnie powinno tutaj być 
zachowane streszczenie na max pół strony o całości projektu}
\end{abstract}

\section*{Zasady i rozwój raportu}
\wskazowka{To jeden dokument rozwijany iteracyjnie. W kolejnych
wersjach poprawiamy również wcześniejsze części i wnioski jeśli będzie taka potrzeba.}


\wskazowka{Raport przygotowujemy w \LaTeX{} (lub jak ktoś nie lubi moze byc doc) i oddajemy jako PDF.
Jedna osoba przekazuje na Delcie raport oraz kod projektu w imieniu
zespołu. Można pracować wspólnie w Overleaf, ale wybór edytora jest dowolny.
Kod powinien mieć instrukcję uruchomienia i opis środowiska (Docker?)}

\wskazowka{Wskazówki, moje uwagi i pola w nawiasach zastępuje własna treść, proszę oczywiście pousuwać to co jest obecnie wypełnione. Szablon jest poglądowy, mogą Państwo zmieniać wedle uznania}

\section{Problem, motywacja i cel}
\subsection{Pytanie predykcyjne}
\wskazowka{Zdefiniuj jednostkę obserwacji, cel i moment predykcji:
dla jakiego obiektu, kiedy i na podstawie jakich dostępnych informacji
przewidujemy jaką wielkość? Jaki jest postawiony nasz cel "biznesowy" ?
Opisz zastosowanie i odbiorce projektu.}

\noindent\textbf{Typ zadania:} \TypZadania.

\noindent\textbf{Kontrakt predykcji:}
\pole{jednostka -- moment -- dostępne cechy -- cel -- zastosowanie etc etc}.

\subsection{Kryterium użyteczności}
\wskazowka{Wyjaśnij, co model ma wnieść ponad prosty punkt odniesienia.}

\section{Dane i analiza wstępna}
\subsection{Źródło i znaczenie kolumn}
\wskazowka{Podaj źródło, autora, odnośnik lub DOI, wersję/datę pobrania
i warunki wykorzystania danych. Opisz pozyskanie, populację,
liczbę rekordów i cech oraz ich znaczenie. Potwierdź zgodność
z wymaganiami zaakceptowanymi przez prowadzącego.}

\begin{table}[H]
\centering
\caption{Skrócony słownik danych: znaczenie, jednostka i dostępność.}
\label{tab:dane}
\footnotesize
\begin{tabularx}{\columnwidth}{@{}lYY@{}}
\toprule
Kolumna & Znaczenie / jednostka & Rola / dostępność\\
\midrule
\texttt{cecha\_1} & \pole{opis} & \pole{przed predykcją?}\\
\texttt{cecha\_2} & \pole{opis} & \pole{przed predykcją?}\\
\texttt{cel} & \pole{etykieta lub jednostka} & Zmienna zależna\\
\bottomrule
\end{tabularx}
\end{table}

\subsection{Jakość i struktura danych}
\wskazowka{Opisz braki, błędy, duplikaty, rzadkie kategorie, wartości
odstające i zależności wierszy. Oczywiście na bazie wykresów. Kazdy wykres musi mieć referencje w tekście oraz opis. Nie wrzucamy kodu do raportu. Klasyfikacja: liczebności i udziały klas.
Regresja: rozkład celu i jednostka.}

\wskazowka{Jeśli braki lub błędy dodano ręcznie, podaj mechanizm,
kolumny, nasilenie, ziarno i pule objęte zmianą. Zachowaj dane
oryginalne. }

\subsection{Najważniejsze obserwacje z EDA}
\wskazowka{Wybierz wykresy wspierające argumentację i zinterpretuj
je w tekście. Analizy służące wyborowi procedury wykonuj na train.
Nie analizuj końcowego testu w celu doboru metod.
Oddziel obserwację od hipotezy i zależność od przyczynowości.}

\section{Przegląd literatury}
\wskazowka{Omów prace dotyczące problemu, danych i adekwatnych metod.
Wskaż zadanie, dane, protokół, wynik i związek z własnym projektem.
Liczby z innych populacji i podziałów nie są bezpośrednio porównywalne.
Dokumentacja biblioteki nie zastępuje badań z publikacji. Najlepiej publikacje operujące na tym samym zbiorze, ale jeśli nie ma to cos podobnego}

\section{Protokół eksperymentu}
\subsection{Podział i granica dopasowania}
\wskazowka{Uzasadnij podział losowy, stratyfikowany, grupowy lub
czasowy względem przyszłego użycia. Podaj liczebności, proporcje,
ziarna i sposób utrwalenia indeksów. Opisz kontrolę powtórzeń
i przenikania informacji. Wiele wierszy nie oznacza automatycznie
wielu niezależnych obiektów. Mocno uzasadnij dlaczego podział jest poprawny i co zostało wzięte pod uwagę}

\begin{itemize}
\item \textbf{Train:} dopasowanie i wewnętrzna CV.
\item \textbf{Validation:} rozwój i wybór rozwiązania; nie jest końcowym potwierdzeniem.
\item \textbf{Test:} ocena zamrożonej procedury.
\end{itemize}

\wskazowka{
Preprocessing, selekcja i uczone transformacje muszą być dopasowane
na właściwej części treningowej, także w CV. Proporcje z wykładów
są przykładem, a nie uniwersalną regułą - można sobie dobierać parametry zgodnie z własną intuicją}

\subsection{Metryka główna i miary pomocnicze}
\wskazowka{Wybór miar i uzasadnienie }

\paragraph{Wariant: klasyfikacja.}
\wskazowka{Określ klasę pozytywną i znaczenie FP/FN. Uzasadnij
miarę decyzji (np. F$_\beta$, recall przy ograniczeniu precision,
balanced accuracy), rankingu (AP, ROC AUC) lub prawdopodobieństw
(log loss, Brier). Samo accuracy może ukrywać ignorowanie
rzadkiej klasy. W zadaniu wieloklasowym podaj wyniki klas i sposób uśredniania.}

\paragraph{Wariant: regresja.}
\wskazowka{Uzasadnij MAE, RMSE lub inne kryterium i podaj jednostkę.
MAE/RMSE mają jednostkę celu, MSE jednostkę kwadratową, R$^2$ może być
ujemne.}

\noindent\textbf{Miara główna:} \pole{nazwa, kierunek, uzasadnienie}.

\noindent\textbf{Miary pomocnicze:} \pole{nazwy i rola}.

\section{Preprocessing i baseline}
\wskazowka{Opisz rzeczywiście użyte imputacje, kodowanie, skalowanie,
filtrowanie lub transformację. Podaj kolejność i miejsce
dopasowania każdego etapu oraz uzasadnienie dla wybranych modeli.
Nie usuwaj trudnych przykładów z validation/test w celu poprawy wyniku ani innych rekordów bez jasnego uzasadnienia}

\wskazowka{Zdefiniuj prosty baseline}

\FloatBarrier
\clearpage
\bibliographystyle{plain}
\bibliography{bibliography}
\end{document}
